
There is also an intrinsic classification with several nonzero coefficient
fibers and a support multiplication different from meet. For a finite group
$\Gamma$, the ordinary monoid semiring $R_0[\Gamma]$ has support semiring
$\mathcal P(\Gamma)$ with union and setwise multiplication, and group $R^A$
over $A\subseteq\Gamma$. Group symbols commute with coefficients, including
when $R$ is noncommutative.

\begin{theorem}[finite group monoid semirings]
\label{thm:finite-group-projectives}
Let $R$ be any unital ring and $\Gamma$ a finite group. Finitely generated
projective left $R_0[\Gamma]$-semimodules are classified by finite
$\Gamma$-posets $P$ and, for one representative $p$ of each orbit, a
finitely generated projective left $R[\Gamma_p]$-module $V_p$, where
$\Gamma_p$ is the stabilizer of $p$. Transport these modules along their
orbits. The support and its groups are
\[
 L=\Id(P),\qquad M_I=\bigoplus_{p\in I}V_p,
\]
with inclusion transitions and the transported $\Gamma$-action. Zero
labels are allowed. Isomorphism means a $\Gamma$-equivariant poset
isomorphism with compatible label isomorphisms; changing an orbit
representative conjugates its stabilizer and transports its module.
No invertibility assumption on $|\Gamma|$ is required.
\end{theorem}

For $R=0$ this says that every finite $\Gamma$-poset gives a projective
$\B[\Gamma]$-semimodule. At a fixed point, the label must be projective
over $R[\Gamma]$. Ordinary projectivity over $R$ alone is insufficient.
Theorem~\ref{thm:finite-group-projectives} evaluates the general support
obstruction for this class; it does not classify arbitrary pairs $(Q,F)$.
